\documentclass[11pt,a4paper]{article}
\usepackage[margin=20mm,headheight=14pt]{geometry}
\usepackage[T1]{fontenc}
\usepackage{helvet}
\renewcommand{\familydefault}{\sfdefault}
\usepackage{amsmath,array,booktabs,tabularx,longtable,xcolor,fancyhdr,hyperref}
\definecolor{navy}{RGB}{14,33,58}
\definecolor{accent}{RGB}{23,140,165}
\definecolor{blue}{RGB}{32,91,151}
\definecolor{soft}{RGB}{239,246,249}
\definecolor{amber}{RGB}{248,239,210}
\hypersetup{colorlinks=true,linkcolor=blue,urlcolor=blue}
\pagestyle{fancy}\fancyhf{}
\lhead{\small Assignment 3: Ducting Layout and Pressure Drop}
\rhead{\small Preliminary design}
\cfoot{\small\thepage}
\newcommand{\nodebox}[2]{\fcolorbox{accent}{soft}{\begin{minipage}[c][11mm][c]{24mm}\centering\footnotesize\textbf{#1}\\[-1mm]\scriptsize #2\end{minipage}}}
\newenvironment{assumptionbox}{\begin{center}\begin{minipage}{0.94\textwidth}\colorbox{amber}{\begin{minipage}{0.95\textwidth}\small\textbf{Assumption boundary.} }{\end{minipage}\end{minipage}\end{center}}
\setlength{\parindent}{0pt}\setlength{\parskip}{5pt}
\renewcommand{\arraystretch}{1.15}
\begin{document}
\begin{titlepage}
\pagecolor{navy}\color{white}
\vspace*{22mm}
{\color{accent}\Large\bfseries ASSIGNMENT 3\par}
\vspace{8mm}
{\Huge\bfseries DUCTING LAYOUT\par}
\vspace{2mm}
{\Huge\bfseries AND PRESSURE-DROP\par}
\vspace{2mm}
{\Huge\bfseries CALCULATION\par}
\vspace{12mm}
{\large Industrial dust-extraction design based on the supplied plant drawings\par}
\vspace{18mm}
\fcolorbox{accent}{navy}{\begin{minipage}{0.87\textwidth}\color{white}\large
Three representative systems \quad|\quad Darcy--Colebrook method\\[3mm]
Critical-path loss \quad|\quad Branch balancing \quad|\quad Fan and motor duty
\end{minipage}}
\vfill
\begin{tabular}{p{48mm}p{95mm}}
Name: & \hrulefill\\[7mm]
Roll No.: & \hrulefill\\[7mm]
Course / Instructor: & \hrulefill\\[7mm]
Submission date: & \hrulefill
\end{tabular}
\vspace{12mm}

{\small Prepared September 2026}
\end{titlepage}
\nopagecolor\color{black}

\section*{Scope statement}
The supplied drawings are low-resolution plant layouts. They establish equipment order and approximate topology, but many dimensions and labels are unreadable. This report therefore presents a transparent preliminary design using stated industrial assumptions. It is suitable as an academic assignment and design basis. A field survey is required before fabrication or purchase.

\section{Design summary}
\begin{center}\small
\begin{tabular}{clrrrrr}
\toprule
System & Duty & Flow & Critical & Required & Selected & Motor\\
 & & m$^3$/h & path & FTP (Pa) & FTP (Pa) & kW\\
\midrule
A & Conveyor extraction & 10000 & A3 & 3402 & 3750 & 18.5 \\
B & Multi-level extraction & 15000 & B1 & 3380 & 3750 & 30.0 \\
C & Cooler and bin extraction & 44000 & C2 & 4690 & 5200 & 110.0 \\
\bottomrule
\end{tabular}
\end{center}
System A is governed by branch A3, System B by B1, and System C by the new-cooler branch G2. FTP denotes fan total pressure at the design flow. Fan static pressure is estimated separately by subtracting the stack outlet velocity pressure.

\section{Drawing interpretation and calculation plan}
\begin{tabularx}{\textwidth}{l X X}
\toprule
Drawing & Readable process intent & Representative model\\
\midrule
1 & Conveyor hoods, collecting duct, bag collector, fan and stack & Four equal pickups in two submains; 10,000 m$^3$/h\\
2 & Multi-level pickups and risers feeding collector, fan and stack & Four elevated pickups in two submains; 15,000 m$^3$/h\\
3 & Existing/new cooler and bin zones, collector/fan and old-fan tie-in & Six equivalent zones; existing/new submains; 44,000 m$^3$/h\\
\bottomrule
\end{tabularx}

The calculation sequence is: define simultaneous hood flows; choose diameters; compute velocity and Reynolds number; solve Colebrook for Darcy friction factor; add friction and fitting losses; compare complete source-to-stack routes; balance the shorter routes; add pressure margin; and calculate fan and motor duty.

\clearpage
\section{Design basis and assumptions}
\begin{tabularx}{\textwidth}{>{\bfseries}l l X}
\toprule
Item & Adopted value & Use\\
\midrule
Air density & 1.20 kg/m$^3$ & Velocity pressure and Reynolds number\\
Dynamic viscosity & $1.81\times10^{-5}$ Pa\,s & Reynolds number\\
Steel roughness & 0.15 mm & Colebrook friction factor\\
Dusty transport velocity & $\geq18$ m/s & Branch, submain and dirty-main sizing\\
Collector differential pressure & 1200 Pa & Base filter allowance\\
Fan total efficiency & 65\% & Shaft power\\
Motor efficiency & 90\% & Electrical input\\
Operation & All pickups simultaneous & No diversity, leakage or buoyancy credit\\
Geometry & Circular internal diameter & Area and hydraulic diameter\\
System C old-fan tie-in & Isolated & No reverse or parallel flow\\
\bottomrule
\end{tabularx}

\subsection{Local-loss coefficients}
\begin{tabularx}{\textwidth}{X r X}
\toprule
Component & $K$ & Reference boundary\\
\midrule
Simple hood entry, Systems A and B & 0.50 & Own branch velocity\\
Equivalent collection zone, System C & 2.00 & Own branch velocity; includes unresolved pickup resistance\\
Long-radius $90^\circ$ bend & 0.25 each & Velocity in its section\\
Incoming combining junction & 0.50 & Incoming-section velocity\\
External gradual adapter & 0.15 & Assigned section velocity\\
Stack terminal including exit & 1.50 & Stack velocity\\
\bottomrule
\end{tabularx}
These are preliminary assumptions, not certified coefficients for the blurred fittings. Exact values depend on geometry and the chosen reference velocity.

\section{Governing equations}
\[
Q=\frac{Q_h}{3600},\qquad A=\frac{\pi D^2}{4},\qquad V=\frac{Q}{A}
\]
\[
Re=\frac{\rho V D}{\mu},\qquad p_v=\frac{\rho V^2}{2}
\]
\[
\frac{1}{\sqrt f}=-2\log_{10}\left(\frac{\epsilon}{3.7D}+\frac{2.51}{Re\sqrt f}\right)
\]
\[
\Delta p_{\rm section}=\left[f\frac{L}{D}+\sum K\right]p_v
\]
Series losses are added. Parallel branch losses are not added together. Each branch route is compared from its pickup to the same stack terminal.
\[
\Delta p_{\rm FTP,req}=\max(\Delta p_{\rm route}),\qquad
K_{\rm bal}=\frac{\Delta p_{\rm crit}-\Delta p_{\rm route}}{p_{v,\rm branch}}
\]
The selected FTP is 10\% above the required value and rounded upward to 50 Pa.
\[
P_{\rm shaft}=\frac{Q\,\Delta p_{\rm FTP}}{1000\,\eta_f},\qquad
P_{\rm input}=\frac{P_{\rm shaft}}{\eta_m}
\]


\clearpage
\section{System A: Conveyor extraction}
\subsection{Interpreted duct layout}

\begin{center}
\begin{tabular}{ccccc}
\nodebox{A1}{2500 m$^3$/h} & $\searrow$ & \nodebox{AS1}{5000 m$^3$/h} & $\searrow$ & \nodebox{Main}{10000 m$^3$/h}\\[2mm]
\nodebox{A2}{2500 m$^3$/h} & $\nearrow$ & & & $\downarrow$\\[4mm]
\nodebox{A3}{2500 m$^3$/h} & $\searrow$ & \nodebox{AS2}{5000 m$^3$/h} & $\nearrow$ & \nodebox{Collector}{1200 Pa}\\[2mm]
\nodebox{A4}{2500 m$^3$/h} & $\nearrow$ & & & $\downarrow$\\
\end{tabular}

\vspace{3mm}
\nodebox{Clean duct}{450 mm}
$\longrightarrow$
\nodebox{Fan}{3750 Pa}
$\longrightarrow$
\nodebox{Stack}{400 mm}
\end{center}
\begin{assumptionbox}
The diagram represents the calculation topology inferred from the scan. It is not an as-built reproduction. Straight lengths, diameters, bend counts and flow allocations are declared design assumptions.
\end{assumptionbox}

\subsection{Section schedule}
\begin{center}\scriptsize
\begin{tabular}{lrrrrrrrrr}
\toprule
ID & $Q_h$ & $D$ & $L$ & $V$ & $f$ & $K$ & $\Delta p_f$ & $\Delta p_K$ & $\Delta p$\\
 & m$^3$/h & mm & m & m/s & & & Pa & Pa & Pa\\
\midrule
A1 & 2500 & 200 & 8.0 & 22.10 & 0.0195 & 1.50 & 229.0 & 439.8 & 668.7 \\
A2 & 2500 & 200 & 6.0 & 22.10 & 0.0195 & 1.25 & 171.7 & 366.5 & 538.2 \\
A3 & 2500 & 200 & 10.0 & 22.10 & 0.0195 & 1.75 & 286.2 & 513.1 & 799.2 \\
A4 & 2500 & 200 & 5.0 & 22.10 & 0.0195 & 1.25 & 143.1 & 366.5 & 509.6 \\
AS1 & 5000 & 300 & 12.0 & 19.65 & 0.0179 & 0.75 & 165.9 & 173.7 & 339.6 \\
AS2 & 5000 & 300 & 6.0 & 19.65 & 0.0179 & 0.75 & 82.9 & 173.7 & 256.7 \\
AM & 10000 & 400 & 10.0 & 22.10 & 0.0167 & 0.65 & 122.3 & 190.6 & 312.8 \\
ACL & 10000 & 450 & 4.0 & 17.47 & 0.0165 & 0.40 & 26.8 & 73.2 & 100.0 \\
AE & 10000 & 400 & 12.0 & 22.10 & 0.0167 & 2.00 & 146.7 & 586.3 & 733.1 \\
\bottomrule
\end{tabular}
\end{center}
All dusty sections lie between \textbf{NaN} and \textbf{NaN m/s}; hence the adopted 18 m/s minimum is satisfied.

\clearpage
\subsection{Complete route comparison and balancing}
\begin{center}\small
\begin{tabular}{llrrr}
\toprule
Source & Route & Total (Pa) & Add (Pa) & $K_{\rm bal}$\\
\midrule
A1 & A1 -- AS1 -- AM -- FILTER -- ACL -- AE & 3354.2 & 47.6 & 0.16 \\
A2 & A2 -- AS1 -- AM -- FILTER -- ACL -- AE & 3223.7 & 178.1 & 0.61 \\
A3 & A3 -- AS2 -- AM -- FILTER -- ACL -- AE & 3401.8 & 0.0 & 0.00 \\
A4 & A4 -- AS2 -- AM -- FILTER -- ACL -- AE & 3112.2 & 289.7 & 0.99 \\
\bottomrule
\end{tabular}
\end{center}
The governing route is
\[
\boxed{\text{A3 -- AS2 -- AM -- FILTER -- ACL -- AE}}
\]
and its calculated resistance is \textbf{3401.8 Pa}. The balancing pressure in the table is $\Delta p_{\rm crit}-\Delta p_{\rm route}$ at the prescribed branch flow.

\subsection{Fan duty}
\begin{center}
\begin{tabularx}{0.82\textwidth}{>{\bfseries}l X}
\toprule
Quantity & Value\\
\midrule
Volume flow & 2.778 m$^3$/s (10000 m$^3$/h)\\
Required fan total pressure & 3401.8 Pa\\
Selected fan total pressure & 3750 Pa\\
Estimated fan static pressure & 3456.8 Pa\\
Shaft power at 65\% efficiency & 16.03 kW\\
Electrical input at 90\% motor efficiency & 17.81 kW\\
Recommended nominal motor & 18.5 kW\\
\bottomrule
\end{tabularx}
\end{center}


\clearpage
\section{System B: Multi-level extraction}
\subsection{Interpreted duct layout}

\begin{center}
\begin{tabular}{ccccc}
\nodebox{B1}{3750 m$^3$/h} & $\searrow$ & \nodebox{BS1}{7500 m$^3$/h} & $\searrow$ & \nodebox{Main}{15000 m$^3$/h}\\[2mm]
\nodebox{B2}{3750 m$^3$/h} & $\nearrow$ & & & $\downarrow$\\[4mm]
\nodebox{B3}{3750 m$^3$/h} & $\searrow$ & \nodebox{BS2}{7500 m$^3$/h} & $\nearrow$ & \nodebox{Collector}{1200 Pa}\\[2mm]
\nodebox{B4}{3750 m$^3$/h} & $\nearrow$ & & & $\downarrow$\\
\end{tabular}

\vspace{3mm}
\nodebox{Clean duct}{560 mm}
$\longrightarrow$
\nodebox{Fan}{3750 Pa}
$\longrightarrow$
\nodebox{Stack}{500 mm}
\end{center}
\begin{assumptionbox}
The diagram represents the calculation topology inferred from the scan. It is not an as-built reproduction. Straight lengths, diameters, bend counts and flow allocations are declared design assumptions.
\end{assumptionbox}

\subsection{Section schedule}
\begin{center}\scriptsize
\begin{tabular}{lrrrrrrrrr}
\toprule
ID & $Q_h$ & $D$ & $L$ & $V$ & $f$ & $K$ & $\Delta p_f$ & $\Delta p_K$ & $\Delta p$\\
 & m$^3$/h & mm & m & m/s & & & Pa & Pa & Pa\\
\midrule
B1 & 3750 & 250 & 16.0 & 21.22 & 0.0186 & 1.75 & 321.2 & 472.8 & 794.1 \\
B2 & 3750 & 250 & 12.0 & 21.22 & 0.0186 & 1.50 & 240.9 & 405.3 & 646.2 \\
B3 & 3750 & 250 & 10.0 & 21.22 & 0.0186 & 1.50 & 200.8 & 405.3 & 606.0 \\
B4 & 3750 & 250 & 8.0 & 21.22 & 0.0186 & 1.25 & 160.6 & 337.7 & 498.3 \\
BS1 & 7500 & 350 & 8.0 & 21.65 & 0.0172 & 1.00 & 110.6 & 281.3 & 391.9 \\
BS2 & 7500 & 350 & 6.0 & 21.65 & 0.0172 & 0.75 & 82.9 & 211.0 & 293.9 \\
BM & 15000 & 500 & 8.0 & 21.22 & 0.0159 & 0.65 & 68.9 & 175.6 & 244.5 \\
BCL & 15000 & 560 & 5.0 & 16.92 & 0.0158 & 0.40 & 24.2 & 68.7 & 92.9 \\
BE & 15000 & 500 & 13.5 & 21.22 & 0.0159 & 2.00 & 116.2 & 540.4 & 656.6 \\
\bottomrule
\end{tabular}
\end{center}
All dusty sections lie between \textbf{NaN} and \textbf{NaN m/s}; hence the adopted 18 m/s minimum is satisfied.

\clearpage
\subsection{Complete route comparison and balancing}
\begin{center}\small
\begin{tabular}{llrrr}
\toprule
Source & Route & Total (Pa) & Add (Pa) & $K_{\rm bal}$\\
\midrule
B1 & B1 -- BS1 -- BM -- FILTER -- BCL -- BE & 3379.9 & 0.0 & 0.00 \\
B2 & B2 -- BS1 -- BM -- FILTER -- BCL -- BE & 3232.1 & 147.9 & 0.55 \\
B3 & B3 -- BS2 -- BM -- FILTER -- BCL -- BE & 3094.0 & 286.0 & 1.06 \\
B4 & B4 -- BS2 -- BM -- FILTER -- BCL -- BE & 2986.3 & 393.7 & 1.46 \\
\bottomrule
\end{tabular}
\end{center}
The governing route is
\[
\boxed{\text{B1 -- BS1 -- BM -- FILTER -- BCL -- BE}}
\]
and its calculated resistance is \textbf{3379.9 Pa}. The balancing pressure in the table is $\Delta p_{\rm crit}-\Delta p_{\rm route}$ at the prescribed branch flow.

\subsection{Fan duty}
\begin{center}
\begin{tabularx}{0.82\textwidth}{>{\bfseries}l X}
\toprule
Quantity & Value\\
\midrule
Volume flow & 4.167 m$^3$/s (15000 m$^3$/h)\\
Required fan total pressure & 3379.9 Pa\\
Selected fan total pressure & 3750 Pa\\
Estimated fan static pressure & 3479.8 Pa\\
Shaft power at 65\% efficiency & 24.04 kW\\
Electrical input at 90\% motor efficiency & 26.71 kW\\
Recommended nominal motor & 30.0 kW\\
\bottomrule
\end{tabularx}
\end{center}


\clearpage
\section{System C: Cooler and bin extraction}
\subsection{Interpreted duct layout}

\begin{center}
\begin{tabular}{ccccc}
\nodebox{G1, G3, G5}{Existing zones} & $\longrightarrow$ & \nodebox{CS1}{22000 m$^3$/h} & $\searrow$ & \nodebox{Main}{44000 m$^3$/h}\\[5mm]
\nodebox{G2, G4, G6}{New zones} & $\longrightarrow$ & \nodebox{CS2}{22000 m$^3$/h} & $\nearrow$ & \nodebox{Collector}{1200 Pa}
\end{tabular}

\vspace{4mm}
\nodebox{Clean duct}{900 mm}
$\longrightarrow$
\nodebox{Fan}{5200 Pa}
$\longrightarrow$
\nodebox{Stack}{850 mm}
\end{center}
\begin{assumptionbox}
The diagram represents the calculation topology inferred from the scan. It is not an as-built reproduction. Straight lengths, diameters, bend counts and flow allocations are declared design assumptions.
\end{assumptionbox}

\subsection{Section schedule}
\begin{center}\scriptsize
\begin{tabular}{lrrrrrrrrr}
\toprule
ID & $Q_h$ & $D$ & $L$ & $V$ & $f$ & $K$ & $\Delta p_f$ & $\Delta p_K$ & $\Delta p$\\
 & m$^3$/h & mm & m & m/s & & & Pa & Pa & Pa\\
\midrule
C1 & 6000 & 300 & 38.0 & 23.58 & 0.0177 & 3.25 & 748.7 & 1084.1 & 1832.8 \\
C2 & 6000 & 300 & 44.0 & 23.58 & 0.0177 & 3.50 & 866.9 & 1167.5 & 2034.4 \\
C3 & 12000 & 450 & 22.0 & 20.96 & 0.0163 & 3.25 & 210.1 & 856.6 & 1066.7 \\
C4 & 12000 & 450 & 28.0 & 20.96 & 0.0163 & 3.50 & 267.4 & 922.5 & 1189.9 \\
C5 & 4000 & 250 & 18.0 & 22.64 & 0.0185 & 3.00 & 409.7 & 922.2 & 1332.0 \\
C6 & 4000 & 250 & 24.0 & 22.64 & 0.0185 & 3.25 & 546.3 & 999.1 & 1545.4 \\
CS1 & 22000 & 625 & 20.0 & 19.92 & 0.0153 & 1.00 & 116.2 & 238.1 & 354.3 \\
CS2 & 22000 & 625 & 16.0 & 19.92 & 0.0153 & 1.00 & 93.0 & 238.1 & 331.0 \\
CM & 44000 & 850 & 40.0 & 21.54 & 0.0143 & 0.65 & 186.8 & 180.9 & 367.7 \\
CCL & 44000 & 900 & 8.0 & 19.21 & 0.0142 & 0.40 & 27.9 & 88.6 & 116.5 \\
CE & 44000 & 850 & 18.0 & 21.54 & 0.0143 & 2.00 & 84.1 & 556.7 & 640.8 \\
\bottomrule
\end{tabular}
\end{center}
All dusty sections lie between \textbf{NaN} and \textbf{NaN m/s}; hence the adopted 18 m/s minimum is satisfied.

\clearpage
\subsection{Complete route comparison and balancing}
\begin{center}\small
\begin{tabular}{llrrr}
\toprule
Source & Route & Total (Pa) & Add (Pa) & $K_{\rm bal}$\\
\midrule
C1 & C1 -- CS1 -- CM -- FILTER -- CCL -- CE & 4512.1 & 178.4 & 0.53 \\
C2 & C2 -- CS2 -- CM -- FILTER -- CCL -- CE & 4690.4 & 0.0 & 0.00 \\
C3 & C3 -- CS1 -- CM -- FILTER -- CCL -- CE & 3746.0 & 944.5 & 3.58 \\
C4 & C4 -- CS2 -- CM -- FILTER -- CCL -- CE & 3845.9 & 844.5 & 3.20 \\
C5 & C5 -- CS1 -- CM -- FILTER -- CCL -- CE & 4011.2 & 679.2 & 2.21 \\
C6 & C6 -- CS2 -- CM -- FILTER -- CCL -- CE & 4201.4 & 489.0 & 1.59 \\
\bottomrule
\end{tabular}
\end{center}
The governing route is
\[
\boxed{\text{C2 -- CS2 -- CM -- FILTER -- CCL -- CE}}
\]
and its calculated resistance is \textbf{4690.4 Pa}. The balancing pressure in the table is $\Delta p_{\rm crit}-\Delta p_{\rm route}$ at the prescribed branch flow.

\subsection{Fan duty}
\begin{center}
\begin{tabularx}{0.82\textwidth}{>{\bfseries}l X}
\toprule
Quantity & Value\\
\midrule
Volume flow & 12.222 m$^3$/s (44000 m$^3$/h)\\
Required fan total pressure & 4690.4 Pa\\
Selected fan total pressure & 5200 Pa\\
Estimated fan static pressure & 4921.6 Pa\\
Shaft power at 65\% efficiency & 97.78 kW\\
Electrical input at 90\% motor efficiency & 108.64 kW\\
Recommended nominal motor & 110.0 kW\\
\bottomrule
\end{tabularx}
\end{center}


\clearpage
\section{Fully worked branch: A3}
For A3, $Q_h=2500$ m$^3$/h, $D=0.200$ m, $L=10$ m and there are three bends.
\[
Q=\frac{2500}{3600}=0.694444\ {\rm m^3/s}
\]
\[
A=\frac{\pi(0.200)^2}{4}=0.03141593\ {\rm m^2},\qquad
V=\frac{0.694444}{0.03141593}=22.1049\ {\rm m/s}
\]
\[
Re=\frac{(1.20)(22.1049)(0.200)}{1.81\times10^{-5}}=293103
\]
\[
p_v=\frac{(1.20)(22.1049)^2}{2}=293.1747\ {\rm Pa}
\]
With $\epsilon/D=0.00015/0.200=0.000750$, the iterated Colebrook equation gives $f=0.01952358$.
\[
\Delta p_f=(0.01952358)\frac{10}{0.200}(293.1747)=286.191\ {\rm Pa}
\]
\[
K=0.50_{\rm hood}+0.50_{\rm junction}+3(0.25_{\rm bend})=1.75
\]
\[
\Delta p_K=(1.75)(293.1747)=513.056\ {\rm Pa}
\]
\[
\boxed{\Delta p_{A3}=286.191+513.056=799.247\ {\rm Pa}}
\]
The complete critical route is
\[
799.247+256.660+312.828+1200+100.033+733.067
=\boxed{3401.834\ {\rm Pa}}
\]

\clearpage
\section{Complete balancing schedule}
\begin{center}\small
\begin{longtable}{llrrrr}
\toprule
System & Branch & Flow & Unbalanced & Add & $K_{\rm bal}$\\
 & & m$^3$/h & Pa & Pa &\\
\midrule
\endhead
A & A1 & 2500 & 3354.2 & 47.6 & 0.16 \\
A & A2 & 2500 & 3223.7 & 178.1 & 0.61 \\
A & A3 & 2500 & 3401.8 & 0.0 & 0.00 \\
A & A4 & 2500 & 3112.2 & 289.7 & 0.99 \\
B & B1 & 3750 & 3379.9 & 0.0 & 0.00 \\
B & B2 & 3750 & 3232.1 & 147.9 & 0.55 \\
B & B3 & 3750 & 3094.0 & 286.0 & 1.06 \\
B & B4 & 3750 & 2986.3 & 393.7 & 1.46 \\
C & C1 & 6000 & 4512.1 & 178.4 & 0.53 \\
C & C2 & 6000 & 4690.4 & 0.0 & 0.00 \\
C & C3 & 12000 & 3746.0 & 944.5 & 3.58 \\
C & C4 & 12000 & 3845.9 & 844.5 & 3.20 \\
C & C5 & 4000 & 4011.2 & 679.2 & 2.21 \\
C & C6 & 4000 & 4201.4 & 489.0 & 1.59 \\
\bottomrule
\end{longtable}
\end{center}
Commissioning should use measured branch flow rather than damper blade angle alone. Adjust the lowest-resistance branches after the collector reaches its normal operating condition, then repeat the traverse until scheduled flows are obtained.

\section{Collector sensitivity}
\begin{center}
\begin{tabular}{crrr}
\toprule
System & Collector $\Delta p$ & Required FTP & Rounded FTP with 10\%\\
 & Pa & Pa & Pa\\
\midrule
A & 800 & 3001.8 & 3350 \\
A & 1200 & 3401.8 & 3750 \\
A & 1600 & 3801.8 & 4200 \\
B & 800 & 2979.9 & 3300 \\
B & 1200 & 3379.9 & 3750 \\
B & 1600 & 3779.9 & 4200 \\
C & 800 & 4290.4 & 4750 \\
C & 1200 & 4690.4 & 5200 \\
C & 1600 & 5090.4 & 5600 \\
\bottomrule
\end{tabular}
\end{center}
A dirtier collector raises system resistance almost one-for-one. Fan control and cleaning strategy must cover the expected differential-pressure band.

\section{Checks, limitations and conclusion}
\begin{itemize}
\item Every dusty section meets the adopted 18 m/s transport criterion.
\item Flow is conserved at all combining junctions.
\item The maximum prescribed-flow route defines each fan duty.
\item Balancing losses use the correct branch velocity-pressure reference.
\item Fan static and total pressure are reported separately.
\item The 10\% design allowance is applied before upward rounding.
\end{itemize}

The resulting preliminary selections are 3750 Pa with an 18.5 kW motor for System A, 3750 Pa with a 30 kW motor for System B, and 5200 Pa with a 110 kW motor for System C. The calculations are internally consistent with the stated topology and assumptions.

Before fabrication, verify every route length, internal diameter, branch angle, elbow radius, hood type and equipment loss. Obtain the collector clean, normal and dirty differential-pressure band and confirm the final operating point on the supplier fan curve at the actual air density.

\section*{References}
\begin{enumerate}
\item NIOSH, \emph{Handbook for Dust Control in Mining}.
\item McGill AirFlow, \emph{Design Advisory 01: Airflow, Friction and Pressure Fundamentals}.
\item McGill AirFlow, \emph{Design Advisory 02: Critical Path and Balancing}.
\item ASHRAE / SMACNA, \emph{Duct Design Fundamentals}.
\item Figueroa et al., hood entry-loss relationship for local exhaust systems.
\item Donaldson, \emph{Differential Pressure in Dust Collectors}.
\end{enumerate}

\begin{assumptionbox}
This submission intentionally separates observed drawing intent from assumed numerical inputs. Replace the assumed values with field measurements and vendor data before procurement.
\end{assumptionbox}
\end{document}